\clearpage
\thispagestyle{empty}
\pdfbookmark[0]{Zur Bestimmung der Konstanten des Bifilarmagnetometers. Ausgabe und Quellen}{bifilar-constants-1840-sources}
\begin{center}\Large\textbf{Zur Bestimmung der Konstanten des Bifilarmagnetometers}\end{center}

Carl Friedrich Gauss’ Abhandlung erschien in den \textit{Resultaten aus den Beobachtungen des magnetischen Vereins} für 1840 als Artikel I. Die KI-Übertragung aller vierzehn Artikel und der historischen Anmerkung in modernes Deutsch mit KI-erstelltem Satz folgt dem deutschen Text in den \textit{Werken}, Bd.~V (Göttingen, 1867), S.~404--426. Die Formeln, Zahlenwerte und historischen Einheiten folgen der deutschen Ausgabe.

\medskip\noindent
Zur Terminologie der Instrumente siehe die Darstellung des Göttinger Magnetischen Vereins bei der Wiechert’schen Erdbebenwarte Göttingen und den Katalogeintrag zur Spiegelfassung für Magnetstäbe des Deutschen Museums, herausgegeben von Julia Bloemer und Benjamin Mirwald, Nr.~1117, Fassung vom 24. August 2016. Der Museumseintrag beschreibt ein Unifilarinstrument.

\medskip\noindent
\href{https://gdz.sub.uni-goettingen.de/id/PPN236006339}{Gesammelte deutsche Ausgabe (1867)}\par
\noindent\href{https://www.erdbebenwarte.de/goettinger-magnetischer-verein/}{Wiechert’sche Erdbebenwarte, Göttinger Magnetischer Verein}\par
\noindent\href{https://digital.deutsches-museum.de/projekte/gruendungssammlung/detail/1117/}{Deutsches Museum, Katalogeintrag zur Spiegelfassung}
\clearpage
